\documentclass[11pt,letterpaper]{article}

\usepackage[utf8]{inputenc}
\usepackage[T1]{fontenc}
\usepackage{lmodern}
\usepackage[spanish,es-noshorthands,es-tabla,es-nodecimaldot]{babel}
\usepackage[margin=2.5cm]{geometry}
\usepackage{amsmath,amssymb}
\usepackage{graphicx}
\usepackage{subcaption}
\usepackage{booktabs}
\usepackage{float}
\usepackage{xcolor}
\usepackage{listings}
\usepackage[hidelinks]{hyperref}

% ---------- Formato del código ----------
\definecolor{fondo}{RGB}{247,247,247}
\definecolor{comentario}{RGB}{95,125,95}
\definecolor{clave}{RGB}{0,70,150}
\definecolor{cadena}{RGB}{160,60,40}

\lstset{
    language=Python,
    basicstyle=\ttfamily\small,
    keywordstyle=\color{clave}\bfseries,
    commentstyle=\color{comentario}\itshape,
    stringstyle=\color{cadena},
    backgroundcolor=\color{fondo},
    frame=single,
    rulecolor=\color{black!25},
    showstringspaces=false,
    breaklines=true,
    columns=fullflexible,
    keepspaces=true,
    captionpos=b,
    literate={á}{{\'a}}1 {é}{{\'e}}1 {í}{{\'i}}1 {ó}{{\'o}}1 {ú}{{\'u}}1 {ñ}{{\~n}}1
}
\renewcommand{\lstlistingname}{Código}

% ---------- Figura con los tres tamaños de un modelo ----------
% #1 = prefijo del archivo en imagenes/, #2 = nombre del modelo, #3 = etiqueta
\newcommand{\tresgrafos}[3]{%
\begin{figure}[H]
    \centering
    \begin{subfigure}{0.49\textwidth}
        \centering
        \includegraphics[width=\linewidth]{imagenes/#1_n50.png}
        \caption{$n = 50$}
    \end{subfigure}\hfill
    \begin{subfigure}{0.49\textwidth}
        \centering
        \includegraphics[width=\linewidth]{imagenes/#1_n200.png}
        \caption{$n = 200$}
    \end{subfigure}

    \medskip
    \begin{subfigure}{0.75\textwidth}
        \centering
        \includegraphics[width=\linewidth]{imagenes/#1_n500.png}
        \caption{$n = 500$}
    \end{subfigure}
    \caption{Grafos generados con el modelo #2.}
    \label{#3}
\end{figure}
\clearpage
}

\title{Proyecto 1\\Biblioteca de generación y manejo de grafos}
\author{Josue gutierrez Andrade\\[4pt]\small Diseño y Análisis de Algoritmos}
\date{}

\begin{document}

\maketitle

\begin{abstract}
Se presenta una biblioteca orientada a objetos, escrita en Python~3, para describir
y generar grafos. La biblioteca define las clases \texttt{Nodo}, \texttt{Arista} y
\texttt{Grafo}, e implementa seis modelos de generación: malla, Erd\H{o}s--Rényi,
Gilbert, geográfico simple, Barabási--Albert y Dorogovtsev--Mendes. Para cada modelo
se generaron grafos de 50, 200 y 500 nodos, que se exportaron a formato GraphViz y se
visualizaron con Gephi.
\end{abstract}

\tableofcontents
\clearpage

% =====================================================================
\section{Estructura de la biblioteca}

El código está organizado en el paquete \texttt{grafos} y un script de apoyo:

\begin{table}[H]
\centering
\begin{tabular}{ll}
\toprule
Archivo & Contenido \\
\midrule
\texttt{grafos/nodo.py} & Clase \texttt{Nodo} \\
\texttt{grafos/arista.py} & Clase \texttt{Arista} \\
\texttt{grafos/grafo.py} & Clase \texttt{Grafo} \\
\texttt{grafos/generadores.py} & Los seis generadores de grafos \\
\texttt{generar\_gv.py} & Script que produce los 18 archivos \texttt{.gv} \\
\texttt{gv/} & Archivos GraphViz generados \\
\texttt{imagenes/} & Visualizaciones hechas con Gephi \\
\bottomrule
\end{tabular}
\caption{Archivos del proyecto.}
\end{table}

% =====================================================================
\section{Representación de los grafos}

\subsection{Clase \texttt{Nodo}}

Un nodo solo almacena su identificador, \texttt{id\_nodo}. La igualdad y el
\emph{hash} se definen a partir de ese identificador, de modo que dos objetos
\texttt{Nodo} con el mismo identificador se consideran el mismo nodo. El identificador
puede ser cualquier valor \emph{hashable}: los modelos aleatorios usan enteros
$0, 1, \dots, n-1$ y la malla usa tuplas $(i, j)$.

\subsection{Clase \texttt{Arista}}

Una arista guarda referencias a sus nodos \texttt{origen} y \texttt{destino}, una
bandera \texttt{dirigida} y un \texttt{peso} opcional (por omisión \texttt{None}).

\subsection{Clase \texttt{Grafo}}

El grafo se representa con una \textbf{lista de adyacencia basada en conjuntos},
acompañada de una lista de aristas. La clase mantiene cuatro estructuras:

\begin{table}[H]
\centering
\begin{tabular}{lll}
\toprule
Atributo & Tipo & Contenido \\
\midrule
\texttt{nodos} & \texttt{dict} & identificador $\rightarrow$ objeto \texttt{Nodo} \\
\texttt{aristas} & \texttt{list} & objetos \texttt{Arista}, en orden de inserción \\
\texttt{\_vecinos} & \texttt{dict} de \texttt{set} & identificador $\rightarrow$ identificadores adyacentes \\
\texttt{\_grado} & \texttt{dict} & identificador $\rightarrow$ número de aristas incidentes \\
\bottomrule
\end{tabular}
\caption{Atributos de la clase \texttt{Grafo}.}
\end{table}

Usar conjuntos en \texttt{\_vecinos} permite consultar si una arista existe en tiempo
$O(1)$ promedio, lo que resulta esencial para los generadores, que hacen esa consulta
cada vez que intentan agregar una arista. La lista \texttt{aristas} es la única
estructura que conserva el peso y la dirección de cada arista, y es la que se recorre
al exportar el grafo.

\subsection{Métodos de \texttt{Grafo}}

\begin{description}
\item[\texttt{agregar\_nodo(id\_nodo)}] Crea el nodo si no existe, inicializa su
conjunto de vecinos vacío y su grado en cero, y devuelve el objeto \texttt{Nodo}.
Si el nodo ya existía, devuelve el existente sin modificar nada.

\item[\texttt{existe\_arista(id\_nodo1, id\_nodo2)}] Devuelve \texttt{True} si
\texttt{id\_nodo2} está en el conjunto de vecinos de \texttt{id\_nodo1}.

\item[\texttt{agregar\_arista(id\_nodo1, id\_nodo2, peso=None)}] Agrega la arista y
la devuelve. Si la arista es un lazo o ya existe, no hace nada y devuelve
\texttt{None}; por lo tanto, los grafos de la biblioteca son siempre \emph{simples}.
Los generadores aprovechan este valor de retorno para saber si la arista realmente
se creó.

\item[\texttt{vecinos(id\_nodo)}] Devuelve el conjunto de identificadores adyacentes.

\item[\texttt{grado(id\_nodo)}] Devuelve el número de aristas incidentes en el nodo.
En grafos dirigidos cuenta tanto las aristas de entrada como las de salida.

\item[\texttt{es\_conexo()}] Hace un recorrido en profundidad iterativo (con una pila)
desde un nodo cualquiera y verifica que se alcancen todos los nodos. Ignora la
dirección de las aristas, así que en grafos dirigidos evalúa conexidad débil.
Su costo es $O(|V| + |E|)$.

\item[\texttt{guardar\_graphviz(archivo)}] Escribe el grafo en formato GraphViz:
primero todos los nodos y después todas las aristas, usando \texttt{graph} y
\texttt{--} para grafos no dirigidos, o \texttt{digraph} y \texttt{->} para dirigidos.
\end{description}

\begin{lstlisting}[caption={Método \texttt{agregar\_arista} de la clase \texttt{Grafo}.}]
def agregar_arista(self, id_nodo1, id_nodo2, peso=None):
    self.agregar_nodo(id_nodo1)
    self.agregar_nodo(id_nodo2)
    if id_nodo1 == id_nodo2 or self.existe_arista(id_nodo1, id_nodo2):
        return None
    arista = Arista(self.nodos[id_nodo1], self.nodos[id_nodo2], self.dirigido, peso)
    self.aristas.append(arista)
    self._vecinos[id_nodo1].add(id_nodo2)
    if not self.dirigido:
        self._vecinos[id_nodo2].add(id_nodo1)
    self._grado[id_nodo1] += 1
    self._grado[id_nodo2] += 1
    return arista
\end{lstlisting}

En un grafo no dirigido la adyacencia se registra en ambos sentidos, pero la arista
aparece una sola vez en la lista \texttt{aristas}. En un grafo dirigido,
\texttt{\_vecinos[u]} contiene únicamente los sucesores de $u$, por lo que las aristas
$u \rightarrow v$ y $v \rightarrow u$ pueden coexistir.

% =====================================================================
\section{Generadores de grafos}

Todos los generadores reciben un parámetro opcional \texttt{dirigido} (por omisión
\texttt{False}) y devuelven un objeto \texttt{Grafo}. Cinco de los seis modelos son
aleatorios y usan el módulo estándar \texttt{random}; como no se fija una semilla,
cada ejecución produce un grafo distinto.

% ---------------------------------------------------------------------
\subsection{Malla: \texttt{grafo\_malla(m, n)}}

Genera una malla de $m \times n$ nodos. Cada nodo se identifica con la tupla $(i, j)$,
con $0 \le i < m$ y $0 \le j < n$, y se conecta con el nodo $(i+1, j)$ y con el nodo
$(i, j+1)$ cuando estos existen. Es el único modelo determinista.

\begin{lstlisting}[caption={Generador de malla.}]
def grafo_malla(m, n, dirigido=False):
    grafo = Grafo(dirigido)
    for i in range(m):
        for j in range(n):
            grafo.agregar_nodo((i, j))
    for i in range(m):
        for j in range(n):
            if i < m - 1:
                grafo.agregar_arista((i, j), (i + 1, j))
            if j < n - 1:
                grafo.agregar_arista((i, j), (i, j + 1))
    return grafo
\end{lstlisting}

El grafo resultante tiene $mn$ nodos y $n(m-1) + m(n-1) = 2mn - m - n$ aristas, y se
construye en tiempo $\Theta(mn)$.

\paragraph{Ejemplo.} La llamada \texttt{grafo\_malla(2, 3)} produce una malla de
6 nodos y $2 \cdot 6 - 2 - 3 = 7$ aristas. Al guardarla con
\texttt{guardar\_graphviz} se obtiene el siguiente archivo:

\begin{lstlisting}[language={},caption={Archivo \texttt{.gv} de la malla de $2 \times 3$.}]
graph G {
    "(0, 0)";
    "(0, 1)";
    "(0, 2)";
    "(1, 0)";
    "(1, 1)";
    "(1, 2)";
    "(0, 0)" -- "(1, 0)";
    "(0, 0)" -- "(0, 1)";
    "(0, 1)" -- "(1, 1)";
    "(0, 1)" -- "(0, 2)";
    "(0, 2)" -- "(1, 2)";
    "(1, 0)" -- "(1, 1)";
    "(1, 1)" -- "(1, 2)";
}
\end{lstlisting}

% ---------------------------------------------------------------------
\subsection{Erd\H{o}s--Rényi: \texttt{grafo\_erdos\_renyi(n, m)}}

Crea $n$ nodos y elige $m$ aristas distintas uniformemente al azar. En cada iteración
se toma un par de nodos distintos con \texttt{random.sample} y se intenta agregar la
arista; solo se cuenta si \texttt{agregar\_arista} la creó, es decir, si no estaba
repetida.

\begin{lstlisting}[caption={Generador de Erd\H{o}s--Rényi.}]
def grafo_erdos_renyi(n, m, dirigido=False):
    max_m = n * (n - 1) if dirigido else n * (n - 1) // 2
    if m > max_m:
        raise ValueError(f"m={m} excede el máximo de aristas posibles ({max_m}) para n={n}")

    grafo = Grafo(dirigido)
    for i in range(n):
        grafo.agregar_nodo(i)

    aristas_creadas = 0
    while aristas_creadas < m:
        u, v = random.sample(range(n), 2)
        if grafo.agregar_arista(u, v) is not None:
            aristas_creadas += 1
    return grafo
\end{lstlisting}

La validación inicial evita un ciclo infinito: si se pidieran más aristas que las
$M = \binom{n}{2}$ posibles (o $M = n(n-1)$ en el caso dirigido), el ciclo nunca
terminaría. Cuando ya hay $k$ aristas, un intento tiene éxito con probabilidad
$1 - k/M$, de modo que el número esperado de intentos es
\[
    \sum_{k=0}^{m-1} \frac{M}{M-k},
\]
que es cercano a $m$ cuando el grafo es disperso ($m \ll M$) y crece conforme $m$ se
aproxima a $M$.

\paragraph{Ejemplo.} \texttt{grafo\_erdos\_renyi(50, 148)} genera un grafo de 50 nodos
con exactamente 148 aristas.

% ---------------------------------------------------------------------
\subsection{Gilbert: \texttt{grafo\_gilbert(n, p)}}

Crea $n$ nodos y, para cada par posible, agrega la arista con probabilidad $p$ de
forma independiente. En el caso no dirigido cada par $\{i, j\}$ se considera una sola
vez (con $i < j$); en el dirigido se consideran los dos pares ordenados.

\begin{lstlisting}[caption={Generador de Gilbert.}]
def grafo_gilbert(n, p, dirigido=False):
    grafo = Grafo(dirigido)
    for i in range(n):
        grafo.agregar_nodo(i)

    for i in range(n):
        for j in range(n):
            if i == j or (not dirigido and j < i):
                continue
            if random.random() < p:
                grafo.agregar_arista(i, j)
    return grafo
\end{lstlisting}

A diferencia del modelo de Erd\H{o}s--Rényi, el número de aristas no es fijo: es una
variable aleatoria binomial con valor esperado $p\binom{n}{2}$ en el caso no dirigido.
El costo es $\Theta(n^2)$, porque se revisan todos los pares.

\paragraph{Ejemplo.} \texttt{grafo\_gilbert(50, 0.1)} genera un grafo de 50 nodos con
alrededor de $0.1 \cdot \binom{50}{2} = 122.5$ aristas en promedio.

% ---------------------------------------------------------------------
\subsection{Geográfico simple: \texttt{grafo\_geografico(n, r)}}

Asigna a cada nodo una posición $(x, y)$ uniforme en el cuadrado unitario y conecta
dos nodos cuando su distancia euclidiana es menor o igual que $r$.

\begin{lstlisting}[caption={Generador geográfico simple.}]
def grafo_geografico(n, r, dirigido=False):
    grafo = Grafo(dirigido)
    posiciones = [(random.random(), random.random()) for _ in range(n)]
    for i in range(n):
        grafo.agregar_nodo(i)

    for i in range(n):
        for j in range(n):
            if i == j or (not dirigido and j < i):
                continue
            dx = posiciones[i][0] - posiciones[j][0]
            dy = posiciones[i][1] - posiciones[j][1]
            if math.hypot(dx, dy) <= r:
                grafo.agregar_arista(i, j)
    return grafo
\end{lstlisting}

La aleatoriedad está solo en las posiciones; una vez fijadas, las aristas quedan
determinadas. Las posiciones se usan únicamente durante la generación y no se guardan
en el grafo. El costo es $\Theta(n^2)$, por la comparación de todos los pares.

\paragraph{Ejemplo.} \texttt{grafo\_geografico(50, 0.22)} genera un grafo de 50 nodos
en el que se conectan los nodos situados a distancia $0.22$ o menor.

% ---------------------------------------------------------------------
\subsection{Barabási--Albert: \texttt{grafo\_barabasi\_albert(n, d)}}

Se implementa la variante del modelo en la que $d$ es el grado máximo de cada nodo.
Los primeros $d$ nodos se conectan todos entre sí. Después, cada nodo nuevo recorre
los nodos existentes en orden aleatorio y se conecta con cada candidato $v$ con
probabilidad
\[
    P(v) = 1 - \frac{\operatorname{grado}(v)}{d},
\]
deteniéndose cuando el nodo nuevo alcanza $d$ aristas.

\begin{lstlisting}[caption={Generador de Barabási--Albert.}]
def grafo_barabasi_albert(n, d, dirigido=False):
    grafo = Grafo(dirigido)
    for i in range(min(d, n)):
        grafo.agregar_nodo(i)
    for i in range(min(d, n)):
        for j in range(i + 1, min(d, n)):
            grafo.agregar_arista(i, j)

    for nuevo in range(d, n):
        grafo.agregar_nodo(nuevo)
        candidatos = list(range(nuevo))
        random.shuffle(candidatos)
        for objetivo in candidatos:
            if grafo.grado(nuevo) >= d:
                break
            probabilidad = 1 - grafo.grado(objetivo) / d
            if random.random() < probabilidad:
                grafo.agregar_arista(nuevo, objetivo)
    return grafo
\end{lstlisting}

Un nodo que ya tiene grado $d$ tiene probabilidad cero de recibir más aristas, por lo
que ningún nodo supera el grado $d$ y el grafo tiene a lo más $nd/2$ aristas. El orden
aleatorio de los candidatos evita favorecer siempre a los nodos con identificador más
bajo. Cada candidato se prueba una sola vez, así que un nodo nuevo puede terminar con
menos de $d$ aristas. Como cada nodo nuevo recorre hasta todos los anteriores, el
costo es $O(n^2)$.

\paragraph{Ejemplo.} \texttt{grafo\_barabasi\_albert(50, 4)} genera un grafo de
50 nodos con grado máximo 4 y, por lo tanto, a lo más 100 aristas.

% ---------------------------------------------------------------------
\subsection{Dorogovtsev--Mendes: \texttt{grafo\_dorogovtsev\_mendes(n)}}

Comienza con un triángulo formado por los nodos 0, 1 y 2. Cada nodo nuevo elige una
arista existente uniformemente al azar y se conecta con sus dos extremos, con lo que
forma un nuevo triángulo.

\begin{lstlisting}[caption={Generador de Dorogovtsev--Mendes.}]
def grafo_dorogovtsev_mendes(n, dirigido=False):
    grafo = Grafo(dirigido)
    for i in range(3):
        grafo.agregar_nodo(i)
    grafo.agregar_arista(0, 1)
    grafo.agregar_arista(1, 2)
    grafo.agregar_arista(2, 0)

    for nuevo in range(3, n):
        grafo.agregar_nodo(nuevo)
        arista = random.choice(grafo.aristas)
        grafo.agregar_arista(nuevo, arista.origen.id_nodo)
        grafo.agregar_arista(nuevo, arista.destino.id_nodo)
    return grafo
\end{lstlisting}

Cada nodo nuevo aporta exactamente dos aristas, así que el grafo tiene siempre
$3 + 2(n-3) = 2n - 3$ aristas y es conexo por construcción. Elegir una arista al azar
equivale a elegir sus extremos con probabilidad proporcional a su grado, lo que
favorece a los nodos mejor conectados. Como \texttt{random.choice} toma tiempo
constante, el costo total es $\Theta(n)$.

\paragraph{Ejemplo.} \texttt{grafo\_dorogovtsev\_mendes(50)} genera un grafo de
50 nodos y $2 \cdot 50 - 3 = 97$ aristas.

% ---------------------------------------------------------------------
\subsection{Resumen}

\begin{table}[H]
\centering
\begin{tabular}{llll}
\toprule
Modelo & Parámetros & Aristas & Costo \\
\midrule
Malla & $m$, $n$ & $2mn - m - n$ & $\Theta(mn)$ \\
Erd\H{o}s--Rényi & $n$, $m$ & $m$ & $\approx m$ intentos si $m \ll \binom{n}{2}$ \\
Gilbert & $n$, $p$ & $p\binom{n}{2}$ en promedio & $\Theta(n^2)$ \\
Geográfico simple & $n$, $r$ & aleatorio & $\Theta(n^2)$ \\
Barabási--Albert & $n$, $d$ & a lo más $nd/2$ & $O(n^2)$ \\
Dorogovtsev--Mendes & $n$ & $2n - 3$ & $\Theta(n)$ \\
\bottomrule
\end{tabular}
\caption{Resumen de los generadores (caso no dirigido).}
\end{table}

% =====================================================================
\section{Ejemplo de uso}

El siguiente programa genera un grafo con el modelo de Gilbert, consulta algunas de
sus propiedades y lo guarda en formato GraphViz:

\begin{lstlisting}[caption={Uso básico de la biblioteca.}]
from grafos import grafo_gilbert

g = grafo_gilbert(50, 0.1)

print(g)                 # Grafo(nodos=50, aristas=..., dirigido=False)
print(g.grado(0))        # grado del nodo 0
print(g.vecinos(0))      # conjunto de vecinos del nodo 0
print(g.es_conexo())     # True si todos los nodos están conectados

g.guardar_graphviz("mi_grafo.gv")
\end{lstlisting}

% =====================================================================
\section{Generación de los archivos GraphViz}

El script \texttt{generar\_gv.py} genera 18 archivos \texttt{.gv}: tres tamaños
(50, 200 y 500 nodos) por cada uno de los seis modelos. Para que las visualizaciones
muestren una sola componente, cada grafo se vuelve a generar hasta que
\texttt{es\_conexo()} devuelve \texttt{True}, con un máximo de 1000 intentos.

Los parámetros de los modelos aleatorios se eligieron ligeramente por encima del
umbral de conectividad, de modo que el grafo resulte conexo casi siempre:

\begin{table}[H]
\centering
\begin{tabular}{ll}
\toprule
Modelo & Parámetro usado \\
\midrule
Erd\H{o}s--Rényi & $m = \operatorname{round}\!\left(\dfrac{n(\ln n + 2)}{2}\right)$ \\[10pt]
Gilbert & $p = \dfrac{\ln n + 3}{n}$ \\[10pt]
Geográfico simple & $r = \sqrt{\dfrac{\ln n + 4}{\pi n}}$ \\[10pt]
Barabási--Albert & $d = 4$ \\[4pt]
Malla & $5 \times 10$, $10 \times 20$ y $20 \times 25$ \\
\bottomrule
\end{tabular}
\caption{Parámetros usados en \texttt{generar\_gv.py}.}
\end{table}

La tabla~\ref{tab:resultados} muestra el número de aristas de los grafos generados.

\begin{table}[H]
\centering
\begin{tabular}{lrrr}
\toprule
 & \multicolumn{3}{c}{Aristas} \\
\cmidrule(l){2-4}
Modelo & $n = 50$ & $n = 200$ & $n = 500$ \\
\midrule
Erd\H{o}s--Rényi & 148 & 730 & 2054 \\
Gilbert & 178 & 814 & 2241 \\
Geográfico simple & 179 & 817 & 2445 \\
Barabási--Albert & 96 & 395 & 996 \\
Dorogovtsev--Mendes & 97 & 397 & 997 \\
Malla & 85 & 370 & 955 \\
\bottomrule
\end{tabular}
\caption{Número de aristas de los 18 grafos generados.}
\label{tab:resultados}
\end{table}

Los valores coinciden con lo esperado: Erd\H{o}s--Rényi tiene exactamente el número de
aristas solicitado, Dorogovtsev--Mendes tiene $2n - 3$, la malla tiene $2mn - m - n$ y
Barabási--Albert queda apenas por debajo de su cota $nd/2$ (100, 400 y 1000).

% =====================================================================
\section{Visualización de los grafos}

Los archivos \texttt{.gv} se visualizaron con Gephi. Las
figuras~\ref{fig:erdos} a~\ref{fig:malla} muestran los 18 grafos generados.

\tresgrafos{1_Erdos_Rényi}{Erd\H{o}s--Rényi}{fig:erdos}
\tresgrafos{2_Gilbert}{Gilbert}{fig:gilbert}
\tresgrafos{3_Geográfico}{geográfico simple}{fig:geografico}
\tresgrafos{4_Barabasi_albert}{Barabási--Albert}{fig:barabasi}
\tresgrafos{5_Dorogovtsev_mendes}{Dorogovtsev--Mendes}{fig:dorogovtsev}
\tresgrafos{6_Malla}{malla}{fig:malla}

\end{document}
